\scanpage{279}
\margin{\unsure{4}}
% archive foliation: 278

Math. (Laboratoire, Ecole) est décidée par l'Assemblée du département de Math.

2°/ La nomination des assistants et Maître-Assistants et la promotion des assistants au poste de Maître-\unsure{Assistant} peut être décidée le cas échéant par une Commission désignée par l'Assemblée. Cette Commission est formée d'enseignants et de représentants d'étudiants du département de Math. Ces derniers comprendront au moins un tiers et au plus une moitié des membres de la Commission et, il est entendu que le plus grand compte sera tenu de l'avis exprimé par ces représentants. La nomination ou la promotion est décidée sur avis favorable d'au moins deux tiers des membres de la Commission.

3°/ La nomination des maîtres de conférences ou des professeurs se fait par l'Assemblée réunie en session plénière. Elle est décidée sur avis favorable de deux tiers au moins des membres votants.

4°/ Dans le cas d'une nomination de Maître de Conférences ou de professeur à l'Ecole de Mathématique, les représentants des étudiants ont un droit de Véto qui devient effectif \unsure{si} deux tiers au moins des représentants d'étudiants se prononcent pour le véto.

5°/ Dans le cas d'une nomination au Laboratoire de Mathématique, et si une majorité des membres de ce laboratoire le demande, la décision est confiée à une Commission désignée par l'Assemblée.
